\documentclass[12pt]{article}
\usepackage[utf8]{inputenc}
\usepackage{amsmath, amssymb, amsfonts}
\usepackage{geometry}
\usepackage{graphicx}
\usepackage{enumitem}
\usepackage{tikz}
\geometry{margin=1in}
\setlist[enumerate]{leftmargin=*}

\title{Quiz 1\\ \vspace{10pt}
\large ECON 320 -- Introduction to Mathematical Economics}
\author{Instructor: Muhebullah Karimzada}
\date{September 17, 2025}

\begin{document}
\maketitle

\section*{Instructions}
\begin{itemize}
    \item Answer all questions. Show all steps clearly to receive full credit.
    \item You may use diagrams where appropriate.
    \item Total time: 15 minutes.
\end{itemize}

\hrulefill


\begin{enumerate}[label=\textbf{Q\arabic*.}]
   
    \item Solve the inequality
    \[
    |2x - 5| < 7
    \]
    and represent the solution set on the real line. Interpret in economic terms if $x$ is quantity and the inequality represents a feasible demand constraint.
   

    \item Consider the feasible set defined by
    \[
    S = \{ (x,y) \in \mathbb{R}^2 \mid x + y \leq 10, \, x \geq 0, \, y \geq 0 \}.
    \]
    \begin{enumerate}[label=(\alph*)]
        \item Sketch the set $S$.  
        \item If $y=2$, what values of $x$ are admissible?  
        \item Interpret this as a resource allocation problem.  
    \end{enumerate}
  


\end{enumerate}




\begin{enumerate}[label=\textbf{Q\arabic*.}, resume]


    \item Consider the demand function given implicitly by
    \[
    Q^d(P,Y) = Y - P^2 = 0,
    \]
    where $P$ is price and $Y$ is income.
    \begin{enumerate}[label=(\alph*)]
        \item Derive $\dfrac{dP}{dY}$.  
        \item What does the sign of this derivative tell us about the relationship between income and equilibrium price?  
    \end{enumerate}


   % \item Let $F(x,p) = 0$ define an equilibrium $x(p)$. Assume $F_x < 0$. 
   % Prove that the sign of the comparative static effect
   % \[
  %  \frac{dx}{dp} = -\frac{F_p}{F_x}
   % \]
  %  depends only on the sign of $F_p$. 
 %   Explain how this allows economists to predict comparative statics without solving explicitly for $x(p)$.
    %\vspace{4cm}


\end{enumerate}

\end{document}
